\documentclass[10pt,letterpaper]{article}
\usepackage[margin=0.68in,top=0.6in,bottom=0.62in]{geometry}
\usepackage{fontspec}
\setmainfont{texgyreheros}[Extension=.otf,UprightFont=*-regular,BoldFont=*-bold,ItalicFont=*-italic,BoldItalicFont=*-bolditalic]
\usepackage{xcolor,enumitem,titlesec,fancyhdr}
\definecolor{accent}{HTML}{176078}
\definecolor{muted}{HTML}{535D65}
\usepackage[unicode,colorlinks=true,urlcolor=accent]{hyperref}
\hypersetup{pdftitle={Pawel Renc — Research CV},pdfauthor={Pawel Renc}}
\setlength{\parindent}{0pt}
\setlength{\parskip}{3pt}
\setlist[itemize]{leftmargin=1.2em,itemsep=2pt,topsep=3pt,parsep=0pt}
\titleformat{\section}{\large\bfseries\color{accent}}{}{0pt}{}[\vspace{2pt}\titlerule]
\titlespacing*{\section}{0pt}{10pt}{6pt}
\pagestyle{fancy}\fancyhf{}
\renewcommand{\headrulewidth}{0pt}
\fancyfoot[L]{\footnotesize\color{muted}Pawel Renc \enspace / \enspace Research CV}
\fancyfoot[R]{\footnotesize\color{muted}\thepage}
\setlength{\footskip}{20pt}
\newcommand{\role}[4]{\textbf{#1}\hfill{\small #2}\\{\color{muted}#3\hfill{\small #4}}\par}
\newcommand{\paper}[3]{\textbf{#1}\\#2\enspace\href{https://doi.org/#3}{DOI: #3}\par\vspace{3pt}}
\begin{document}
\raggedright
{\fontsize{29}{32}\selectfont\bfseries Pawel Renc}\\[4pt]
{\large\color{accent}Healthcare AI \enspace / \enspace Foundation Models \enspace / \enspace ML Systems}\\[6pt]
{\small Boston, MA \enspace · \enspace \href{mailto:rencpawe@gmail.com}{rencpawe@gmail.com}
\enspace · \enspace \href{https://github.com/prenc}{GitHub}
\enspace · \enspace \href{https://www.linkedin.com/in/pawel-renc}{LinkedIn}
\enspace · \enspace \href{https://scholar.google.com/citations?user=Ip7ldAEAAAAJ}{Google Scholar}}
\section{Research profile}
Research fellow at Massachusetts General Hospital working on generative models of longitudinal electronic health records and the infrastructure for training and evaluating them. Co-developer of ETHOS and ARES, combining healthcare-AI research with industry experience at Snowflake and high-performance computing research at ETH Zürich.
\section{Research \& industry experience}
\role{Research Fellow}{May 2026–present}{Massachusetts General Hospital · Sitek Lab}{Boston, MA}
\begin{itemize}
\item Develop EHR foundation models and end-to-end infrastructure for large-scale data processing, training and inference.
\item Maintain the ETHOS research codebase, spanning tokenization, model training, evaluation and synthetic timeline generation.
\end{itemize}
\role{Research Scientist}{Oct 2025–May 2026}{Snowflake · AI Research}{Bellevue, WA}
\begin{itemize}
\item Designed and developed an agentic medical-coding system for assigning diagnoses from clinical notes.
\item Partnered with forward-deployed engineering teams on healthcare applications of medical foundation models.
\end{itemize}
\role{AI Research Intern}{Apr–Oct 2025}{Snowflake}{Bellevue, WA}
\begin{itemize}
\item Developed and fine-tuned EHR foundation models; built training and inference workflows with Transformers, vLLM and Ray.
\end{itemize}
\role{Visiting PhD Researcher}{Mar 2023–Apr 2025}{Massachusetts General Hospital · Harvard Medical School}{Boston, MA}
\begin{itemize}
\item Co-developed ETHOS for generative patient-timeline modeling and zero-shot health prediction; contributed to timeline representation, tokenization, transformer training and inference.
\item Led software and formal analysis for ARES, extending timeline simulation to adaptive risk estimation.
\end{itemize}
\role{Visiting PhD Researcher}{2022–2023}{ETH Zürich · Scalable Parallel Computing Laboratory}{Zürich, Switzerland}
\begin{itemize}
\item Researched GPU execution of graph neural networks and global tensor formulations; coauthored work published at SC 2023.
\end{itemize}
\role{PhD Researcher}{Jan 2022–Feb 2023}{Sano Centre for Computational Medicine}{Kraków, Poland}
\begin{itemize}
\item Developed machine-learning methods for disease-phenotype discovery through patient-data clustering, in collaboration with University of Glasgow clinicians.
\end{itemize}
\section{Technical expertise}
\textbf{Modeling:} EHR foundation models, generative AI, transformers, graph neural networks.\\
\textbf{ML systems:} PyTorch, Transformers, vLLM, Ray; distributed training and inference.\\
\textbf{Engineering:} Python, C/C++, CUDA, Julia, Bash; Polars, SQL, Linux and GPU computing.
\newpage
{\large\bfseries Pawel Renc}\hfill{\small\color{muted}Selected research, education \& service}
\section{Selected publications}
\paper{Foundation model of electronic medical records for adaptive risk estimation}
{First author · \textit{GigaScience}, 2025.}{10.1093/gigascience/giaf107}
\paper{Zero shot health trajectory prediction using transformer}
{First author · \textit{npj Digital Medicine}, 2024.}{10.1038/s41746-024-01235-0}
\paper{High-Performance and Programmable Attentional Graph Neural Networks\\with Global Tensor Formulations}
{Coauthor · \textit{SC}, 2023.}{10.1145/3581784.3607067}
\paper{Towards efficient GPGPU Cellular Automata model implementation\\using persistent active cells}
{Coauthor · \textit{Journal of Computational Science}, 2022.}{10.1016/j.jocs.2021.101538}
\paper{EBIC.JL — an efficient implementation of evolutionary biclustering algorithm in Julia}
{Coauthor · \textit{GECCO Companion}, 2021.}{10.1145/3449726.3463197}
\section{Research software}
\textbf{\href{https://github.com/ipolharvard/ethos-ares}{ETHOS / ARES}} — EHR tokenization, transformer training and health-trajectory inference.
Research software supporting generative patient-timeline modeling.\\
\textbf{\href{https://github.com/EpistasisLab/EBIC.jl}{EBIC.jl}} — GPU-enabled evolutionary biclustering in Julia; developed during master's research.
\section{Education}
\textbf{Doctorate · Engineering and Technology}\hfill May 2026\\
AGH University of Krakow · Discipline: Information and Communication Technology\\
{\small Dissertation: \textit{Foundation model for electronic health records}.\\
Supervisors: Jarosław Wąs and Arkadiusz Sitek. Degree conferred May 21, 2026.}\par
\textbf{Master's degree · Computer Science}\hfill 2021\\
AGH University of Science and Technology · System modelling and data analysis\par
\textbf{Bachelor's degree · Computer Science}\hfill 2020\\
AGH University of Science and Technology
\section{Selected service \& teaching}
\textbf{Peer review:} IEEE Journal of Biomedical and Health Informatics (two rounds, 2025);
IEEE Transactions on Medical Imaging (2023); BMJ Digital Health \& AI (2026).\\
\textbf{Grand rounds:} Applications of Machine Learning in Healthcare, DHR Health, April 2024.\\
\textbf{Research and Teaching Assistant, AGH}\hfill 2020–2021\\
Taught Unix, Operations Research and Discrete Systems Modeling.
\section{Earlier engineering experience}
\textbf{QED Software · Data Scientist}\hfill 2020–2021\\
Built ML methods for detecting erroneous table joins and distinguishing bots from human game players.\\
\textbf{Intel · Software Engineering Intern}\hfill 2019\\
Developed Python automation for integration, unit and system-level testing.\\
\textbf{GrowMagick · Software Engineer}\hfill 2019–2020\\
Built data pipelines and anomaly-detection workflows for an advertising-optimization R\&D project.
\end{document}
